%% ═══════════════════════════════════════════════════════════════════════════
%% FUTURE WORK (FULL VERSION) -- reference copy, not included in main.tex
%% ═══════════════════════════════════════════════════════════════════════════

% The thesis uses a condensed version of this section.
% This file preserves all 10 future work directions for reference.

\subsection{Future Work}\label{sec:future_work_full}

Several directions could extend this framework.

\paragraph{Predictive regime signal}
The current regime signal $\pi_t^{\text{filter}}$ estimates the probability that the market is in panic \emph{now}. A natural extension is to predict the regime one step ahead: $P(s_{t+1} = 1 \mid \mathbf{z}_{1:t})$, which can be obtained directly from the transition matrix as $\sum_k \pi_t^{\text{filter}}(k) \cdot P_{k,1}$. A predictive signal would allow the portfolio to reposition \emph{before} a regime shift rather than reacting to it, potentially reducing turnover and improving timing around transition months.

\paragraph{Adaptive regime estimation}
The current HMM is trained on 1990--2010 data and its parameters are held fixed during the test period. This assumes that the statistical properties of calm and panic regimes remain stationary over time. If the market undergoes structural change, for instance a permanent increase in baseline volatility due to algorithmic trading, higher-frequency information flow, or shifting monetary policy regimes, the trained thresholds may become miscalibrated. Concretely, features that were one standard deviation above the training mean may become the new normal, causing the model to classify structurally calm periods as panic. This concern is not hypothetical: the sustained elevation of $\pi_t^{\text{filter}}$ in the most recent years of the sample may partly reflect such drift rather than genuine prolonged stress. Rolling or expanding-window re-estimation of the HMM, or an online Bayesian updating scheme that allows the emission parameters to evolve over time, would address this limitation and likely improve the reliability of the regime signal in extended out-of-sample periods.

\paragraph{International replication}
The current analysis is limited to US equities. Testing whether the continuation-to-reversal shift and the nonlinearity finding replicate in international markets (e.g., Europe, Japan, emerging markets) would strengthen the generalizability of the mechanism.

\paragraph{Exposure-scaling circuit breaker}
The strategy's primary vulnerability is a prolonged economic deterioration where beaten-down stocks do not recover. A natural extension is to combine XGB's stock selection with a D\&M-style exposure-scaling overlay that reduces overall portfolio exposure when a sustained bear market indicator (e.g., consecutive months of elevated $\pi_t^{\text{filter}}$) is triggered. Preliminary tests on the current sample show no improvement, since the 2011--2025 period contains no prolonged recession and the circuit breaker only cuts profitable panic months; the value of this mechanism would emerge in a sustained downturn outside the current sample.

\paragraph{Long-short symmetry}
In both regimes, the short leg's z-score curve approximately mirrors the long leg's: the model applies symmetric but opposite selection criteria to both sides of the portfolio. It remains an open question whether this symmetry is an inherent property of cross-sectional momentum or an artefact of the decile-based construction.

\paragraph{Regime-conditional feature sets}
The feature-set ablation (Section~\ref{sec:fund_ablation}) shows that adding ten fundamental characteristics to XGB reduces Sharpe from 1.11 to 0.89, and that fundamentals alone produce a calm Sharpe close to the baseline (0.86 vs.\ 0.84) but a panic Sharpe of only 0.38. Firm-level characteristics describe stable attributes that cannot flip direction across regimes, so they dilute the joint model rather than augment it. This suggests that the optimal feature set may itself be regime-dependent: momentum and the regime signal carry the mechanism during panic, where reversal dominates, while fundamentals might add value during calm if used differently, for example as a quality filter applied after the momentum ranking rather than as an input to the model. Exploring regime-conditional feature sets, in which the model uses different inputs depending on the market state, could recover the information fundamentals plausibly carry without diluting the reversal mechanism.

\paragraph{Alternative nonlinear models}
This thesis uses XGBoost as the nonlinear model, but the finding that the regime-momentum interaction requires nonlinear modelling should hold for other flexible model classes. Comparing XGBoost against alternative architectures such as LightGBM and neural networks (including attention-based models) would clarify whether the regime-conditional shift is captured equally well by all nonlinear methods or whether tree-based models have a specific advantage due to their discrete splitting structure. If neural networks capture the same shift, it would strengthen the conclusion that the finding is about the data, not the model.

\paragraph{Sequential regime conditioning}
XGBoost observes only the current month's features and has no memory of past states. Two months with $\pi_t^{\text{filter}} = 0.75$ are treated identically regardless of whether the signal just spiked from 0.1 (a fresh crash, where reversal is most likely) or has been elevated for six consecutive months (prolonged stress, where reversal may fail). In practice these are very different situations. A sequential model (for example, an LSTM or a simpler approach that includes lagged values of $\pi_t^{\text{filter}}$ as additional features) could learn to distinguish fresh stress from sustained deterioration, potentially improving the model's ability to detect when the reversal mechanism is breaking down.

\paragraph{Alternative emission distributions}
The current HMM uses Gaussian emissions, which are robust in practice (Appendix~\ref{app:emission_robustness}). Exploring alternative emission distributions, such as multivariate Student-$t$ or skew-normal specifications, may better accommodate the heavy tails and asymmetry present in financial stress indicators and could improve regime identification in samples with more extreme tail events.

\paragraph{Transaction costs and implementation}
The analysis applies a uniform 10 basis point one-way transaction cost, but does not model the higher costs that may affect smaller or less liquid stocks. In long-short construction, XGB's turnover is higher than in long-only due to regime-conditional changes in stock selection. Incorporating more realistic, stock-specific transaction cost estimates would provide a more accurate picture of implementability.
